%%%%%%%%%%%%%%%%%%%%%%%%%%%%%%%%%%%%%%%%%%%%%%%%%%%%%%%%%%%%%%%%%%%%%%%%%%%%%%%%%%%
%% This project aims to create  a custom  UDL template for presentation.                %%
%% author: BOUKLI HACENE Sofiane - Prof in Computer Science Departement (UDL)   %%
%% contacts:                                                                     %%
%%    e-mail: boukli@univ-sba.dz                                                   %%
 
%%%%%%%%%%%%%%%%%%%%%%%%%%%%%%%%%%%%%%%%%%%%%%%%%%%%%%%%%%%%%%%%%%%%%%%%%%%%%%%%%%%

\documentclass{libs/udl_format}
% Inserting the preamble file with the packages
\input{libs/preamble.tex}
% Inserting the references file
\bibliography{references.bib}

% Title
\title[short title]{\huge\textbf{Titre de la présentation}}
% Subtitle
\subtitle{Sous titre}
% Author of the presentation
\author{Nom Auteur \& Nom Auteur}
% Institute's Name
\institute[UDL]{
    % email for contact
    \normalsize{\email{email@univ-sba.dz}}
    \newline
    \speciality{ Ma spécialité }
    % Department Name
    \department{Département d'informatique}
    \newline
    % university name
    \udl
}
% date of the presentation
% \date{\today}

%%%%%%%
%%%%%%




%%%%%%%%%%%%%%%%%%%%%%%%%%%%%%%%%%%%%%%%%%%%%%%%%%%%%%%%%%%%%%%%%%%%%%%%%%%%%%%%%%
%% Start Document of the Presentation                                           %%               
%%%%%%%%%%%%%%%%%%%%%%%%%%%%%%%%%%%%%%%%%%%%%%%%%%%%%%%%%%%%%%%%%%%%%%%%%%%%%%%%%%
\begin{document}
% insert the code style
\input{libs/code_style}

%% ---------------------------------------------------------------------------
%%%%%$$



%%%%%%%%%



% First frame (with tile, subtitle, ...)
\begin{frame}{}
\centering{\calligra{Mémoire de Master } } %% Thèse de Doctorat
    \maketitle
 
    %%%%%%%%%%%%%%%
  %\hspace{-0.5cm} devant le jury :
  \hspace{1.5cm}  {
 \resizebox{0.8\textwidth}{!}{
  \begin{tabular}{llll}
			Pr.  & \textsc{Nom et prénoms jury} & UDL SBA & Président du jury \\
				Pr.  & \textsc{Nom et prénoms jury} & Université jury & Examinateur \\
			Dr.  & \textsc{Nom et prénoms jury} & UDL SBA & Directeur de thèse  \\
			Pr.  & \textsc{Nom et prénoms jury} & UDL SBA & CO-Directeur de thèse \\
		\end{tabular}   
}
}
 
\normalsize
\end{frame}

%% ---------------------------------------------------------------------------
% Second frame
\begin{frame}{Sommaire}
    \begin{multicols}{2}
        \tableofcontents
    \end{multicols}
\end{frame}

%% ---------------------------------------------------------------------------
% This presentation is separated by sections and subsections
\section{Section1}
\begin{frame}{Explictions}
    % itemize
    Il s'agit d'un modèle qui peut être utilisé :
    \begin{itemize}
        \item Présentation des articles universitaires
        \item Soumission des sujets
        \item Présentations de thèses et de mémoires
    \end{itemize}

    \vspace{0.4cm} % vertical space
    
    % enumeration
    
        Pour utiliser correctement ce modèle, il est important que :
      \begin{enumerate}
        \item Avoir une connaissance minimale de LaTeX
        \item Lire les commentaires dans le modèle (explications)
        \item Lire le fichier README.md (documentation)
    \end{enumerate}

    \vspace{0.2cm}
\example{Ceci est un exemple de texte!} \emph{Emphasis Text!}
    
\end{frame}

%% ---------------------------------------------------------------------------
\subsection{Sous Section1}
\begin{frame}{Créer des blocs}
    % Blocks styles
    \begin{block}{Bloc standard}
        Mettre en bloc le corps du texte.
    \end{block}

    \begin{alertblock}{Bloc Alerte}
        Mettre en bloc le corps du texte.
    \end{alertblock}

    \begin{exampleblock}{Bloc exemple}
         Mettre en bloc le corps du texte.
    \end{exampleblock}   
\end{frame}

%% ---------------------------------------------------------------------------
\subsection{Sous Section 2}
\begin{frame}{Créer des boîtes}
    \successbox{test ou boîte de réussite}

    \pause

    \alertbox{test ou boîte d'alerte}

    \pause

    \simplebox{test ou boîte d'exemple}
\end{frame}

%% ---------------------------------------------------------------------------
\subsection{Sous Section 3}
\begin{frame}{Créer des Algorithmes (Pseudocode)}
    \begin{algorithm}[H]
        \SetAlgoLined
        \LinesNumbered
        \SetKwInOut{Input}{input}
        \SetKwInOut{Output}{output}
        \Input{x: float, y: float}
        \Output{r: float}
        \While{True}{
          r = x + y\;
          \eIf{r >= 30}{
           ``La valeur de $r$ est supérieure ou égale à 30.''\;
           break\;
           }{
           ``La valeur de $r$ = '', r\;
          }
         } 
         \caption{Exemple d'algorithme }
    \end{algorithm}
\end{frame}

%% ---------------------------------------------------------------------------

\begin{frame}{Insertion de Programmes Python }
    \lstset{language=Python}
    \lstinputlisting[language=Python]{code/main.py}
\end{frame}

%% ---------------------------------------------------------------------------
\begin{frame}{Insertion de Programmes C }
    \lstinputlisting[language=C]{code/source.c}
\end{frame}

%% ---------------------------------------------------------------------------
\begin{frame}{Insertion de Programmes JAVA}
    \lstinputlisting[language=Java]{code/helloworld.java}
\end{frame}

%% ---------------------------------------------------------------------------
\begin{frame}{Insertion de Programmes HTML}
    \lstinputlisting[language=HTML]{code/index.html}
\end{frame}

%% ---------------------------------------------------------------------------
% This frame show an example to insert multicolumns
\section{Multicolonnes}
\begin{frame}{Section 2 - Multicolonnes}
    \begin{columns}{}
        \begin{column}{0.5\textwidth}
            \justify
            Vous pouvez placer plus d'une colonne en utilisant le $\backslash$begin\{column\}\{\} e $\backslash$end\{column\}
        \end{column}
        \begin{column}{0.5\textwidth}
            \justify
            Cependant, l'espacement doit être proportionnel entre les colonnes afin que ces dernières ne se chevauchent pas. L'espacement est donné par le deuxième argument de la $\backslash$begin.
        \end{column}
    \end{columns}    
\end{frame}

%% ---------------------------------------------------------------------------
% This frame show an example to insert figures
\section{Images}
\begin{frame}{Section 3 - Figures}
    \begin{figure}
        \centering
        \caption{Logo de L'UDL.}
        \includegraphics[scale=0.3]{libs/emblemudl.pdf}
        \source{Récupéré sur le site officiel de l'UDL \cite{siteudl} \cite{einstein}}
        \label{fig:udl_emblem}
    \end{figure}
\end{frame}

%% ---------------------------------------------------------------------------
% Reference frames
\begin{frame}[allowframebreaks]
    \frametitle{Références}
    \printbibliography
\end{frame}

%% ---------------------------------------------------------------------------
% Final frame
\begin{frame}{}
    \centering
    \huge{\textbf{\example{Merci de votre attention !}}}
    
    \vspace{1cm}
    
    \Large{\textbf{Contact:}}
    \newline
    \vspace*{0.5cm}
    \large{\email{email@univ-sba.dz}}
\end{frame}

\end{document}